\section{Conclusion}

This paper has addressed the critical challenges of signal detection in extremely large-scale multiple-input multiple-output (XL-MIMO) systems operating in the near-field regime. Specifically, conventional full-array processing approaches and fully-connected graph neural network (GNN) detectors face prohibitive computational complexity and severe noise accumulation due to the spatial non-stationarity inherent in massive antenna deployments.

To overcome these limitations, a low-complexity detection framework, termed Graph-VR-Det, was proposed. The methodology leveraged a pilot-driven dynamic visibility region (VR) discovery algorithm to identify localized user-antenna connections, thereby constructing a sparse bipartite graph topology. By restricting the evidence processing of a sparse message passing neural network (Sparse-GNN) exclusively to these valid edges, the framework significantly reduced computational overhead. Furthermore, a frequency-aware feature embedding mechanism was integrated to compensate for the wideband near-field beam splitting effect at terahertz (THz) frequencies, enabling the network to dynamically adjust to spatial shifts across subcarriers.

Extensive numerical evaluations revealed several core findings. The proposed sparse architecture achieved an approximate $3\times$ inference speedup compared to fully-connected baselines, effectively decoupling the detection complexity from the total array size. Notably, at an optimal sparsity level, the selective pruning of weak connections induced a graph regularization effect that filtered out noise from non-visible antennas, allowing the Sparse-GNN to surpass the bit error rate (BER) performance of dense topologies. Additionally, the localized nature of the message passing operations facilitated robust zero-shot scalability, enabling the model to seamlessly generalize to a 1024-antenna system without requiring retraining.

Despite these advantages, the reliance on an energy-based VR discovery mechanism introduces an error floor at high signal-to-noise ratios (SNRs) in wideband scenarios. Future research should focus on mitigating this high-SNR performance gap by developing more sophisticated, channel-resilient edge selection criteria. Furthermore, extending the proposed sparse framework to accommodate multi-antenna users and exploring hardware-efficient implementations of the required sparse scatter operations represent promising directions for practical real-time deployments.